\documentclass[10pt,a4paper]{article}
\usepackage{nexusS5}
\setnexusheader{Entrée en Seconde générale et technologique — Stage de pré-rentrée}{Évaluation finale — Ahmed BAKIR}
\begin{document}
\nexuscover{ÉVALUATION FINALE --- DIAGNOSTIC DE SORTIE}{Durée : 45 minutes}{Ahmed BAKIR}{Entrée en Seconde générale et technologique}{Mathématiques --- à traiter individuellement}
\par\noindent\begin{tabularx}{\linewidth}{>{\bfseries}p{22mm} X >{\bfseries}p{16mm} X}
\toprule Nom et prénom & \rule{62mm}{0.32pt} & Date & \rule{34mm}{0.32pt} \\ \bottomrule
\end{tabularx}

\begin{goalbox}
\textbf{Ce document est un diagnostic de sortie, pas un classement.} Il mesure ce qui est
désormais installé, ce qui reste fragile, et il sert à écrire le plan des quatre premières
semaines de la rentrée.
\end{goalbox}

\begin{methodbox}
\textbf{Consignes}
\begin{itemize}
\item Durée : \textbf{45 minutes}. Partie A : environ 10 min. Partie B : environ 20 min. Partie C : environ 15 min.
\item Travailler seul, sans document et sans les cartes de rappel de la séance.
\item Écrire la relation ou la propriété utilisée avant le calcul.
\item Une question peut être laissée de côté puis reprise ; il vaut mieux la laisser vide que d'écrire au hasard.
\item Trois lignes « Certitude » seulement : \conf\ , sur la même échelle qu'en début de stage.
\end{itemize}
\end{methodbox}

\newpage
\phasehead{Partie A --- Automatismes et connaissances essentielles}{environ 10 min}{Questions courtes. Écrire le résultat, sans rédaction longue.}
\begin{enumerate}
\needspace{22mm}
\item[\textbf{A1.}] Calculer $-4 + 5\times(-3) - (-6)$.
\worklines{2}
\par\vspace{1.2mm}
\needspace{22mm}
\item[\textbf{A2.}] Calculer $\dfrac{3}{8} + \dfrac{5}{12}$.
\worklines{2}
\par\vspace{1.2mm}
\needspace{22mm}
\item[\textbf{A3.}] Simplifier $\dfrac{10^{5}\times 10^{-2}}{10^{4}}$ et donner le résultat sous la forme d'une puissance de $10$.
\worklines{2}
\par\vspace{1.2mm}
\needspace{22mm}
\item[\textbf{A4.}] Développer $(x - 5)^2$.
\worklines{2}
\par\vspace{1.2mm}
\needspace{22mm}
\item[\textbf{A5.}] Calculer $(-3)^3 + 2\times(-4)$.
\worklines{2}
\par\vspace{1.2mm}
\needspace{22mm}
\item[\textbf{A6.}] Écrire $0{,}000205$ en notation scientifique.
\worklines{2}
\par\vspace{1.2mm}
\end{enumerate}
\certbox{Certitude sur l'ensemble de la partie A :}
\needspace{68mm}\par\vspace{3mm}
\phasehead{Partie B --- Application et raisonnement}{environ 20 min}{Écrire la relation, la propriété ou la règle avant le calcul, puis contrôler.}
\begin{enumerate}
\needspace{22mm}
\item[\textbf{B1.}] \begin{enumerate}
\item Résoudre $5x - 8 = 2x + 7$, puis vérifier la solution.
\item Résoudre $(2x + 7)(x - 3) = 0$ en indiquant la propriété utilisée.
\item Vérifier la solution entière de cette seconde équation.
\end{enumerate}
\worklines{5}
\par\vspace{1.2mm}
\needspace{22mm}
\item[\textbf{B2.}] Un article coûte $250$ TND. Son prix augmente de $12\,\%$, puis le nouveau prix baisse de $12\,\%$.
\begin{enumerate}
\item Calculer le prix après chaque étape.
\item Le prix final est-il égal au prix initial ? Justifier à l'aide du coefficient multiplicateur global.
\end{enumerate}
\worklines{5}
\par\vspace{1.2mm}
\needspace{22mm}
\item[\textbf{B3.}] Le triangle $ABC$ est rectangle en $A$, avec $AB = 9$ cm et $AC = 12$ cm.
Le point $M$ appartient à $[AB]$ et le point $N$ à $[AC]$, avec $(MN)$ parallèle à $(BC)$ et $AM = 3$ cm.
\begin{enumerate}
\item Calculer $BC$.
\item Calculer $AN$ et $MN$ en citant le théorème utilisé.
\item En déduire le périmètre du triangle $AMN$.
\end{enumerate}
\worklines{5}
\par\vspace{1.2mm}
\needspace{22mm}
\item[\textbf{B4.}] \begin{enumerate}
\item Développer et réduire $(2x - 3)(x + 5)$.
\item Contrôler le résultat pour $x = 1$, en calculant séparément les deux écritures.
\end{enumerate}
\worklines{6}
\par\vspace{1.2mm}
\end{enumerate}
\certbox{Certitude sur la question B2 :}
\needspace{68mm}\par\vspace{3mm}
\phasehead{Partie C --- Transfert et synthèse}{environ 15 min}{Reconnaître la notion utile, mobiliser plusieurs acquis, conclure par une phrase.}
\begin{enumerate}
\needspace{22mm}
\item[\textbf{C1.}] \textbf{L'écran d'affichage.} Un écran rectangulaire a pour largeur $x$ mètres et pour
hauteur $x - 1$ mètres, avec $x > 1$.
\begin{enumerate}
\item Écrire l'aire $A(x)$ de l'écran, puis la développer et la réduire.
\item Calculer $A(4)$, puis vérifier ce résultat par le calcul direct de l'aire d'un écran de $4$ m sur $3$ m.
\item Calculer la longueur de la diagonale de cet écran de $4$ m sur $3$ m, en citant le théorème utilisé.
\item L'écran est affiché $1\,250$ TND. Une remise de $15\,\%$ est accordée, puis une taxe de $19\,\%$ s'applique
sur le prix remisé. Calculer le prix final, puis dire s'il est supérieur ou inférieur au prix affiché.
Justifier à l'aide du coefficient multiplicateur global.
\end{enumerate}
\worklines{12}
\par\vspace{1.2mm}
\needspace{22mm}
\item[\textbf{C2.}] Un forfait coûte $60$ TND.
\begin{enumerate}
\item Calculer son prix après une hausse de $15\,\%$.
\item Un élève affirme qu'une hausse de $15\,\%$ suivie d'une baisse de $15\,\%$ ramène à $60$ TND.
Réfuter par un calcul.
\end{enumerate}
\worklines{5}
\par\vspace{1.2mm}
\end{enumerate}
\certbox{Certitude sur la question C1 :}
\brouillon{\dimexpr\textheight/3\relax}
\newpage
\sectiontitle{Après l'évaluation --- à remplir en deux minutes}
\begin{methodbox}Ces trois lignes ne sont pas notées. Elles servent à construire le plan des quatre premières semaines de la rentrée.\end{methodbox}
\par\noindent\begin{tabularx}{\linewidth}{>{\bfseries}p{62mm} X}
\toprule
La question que j'ai trouvée la plus sûre & \rule{85mm}{0.32pt} \\
La question sur laquelle j'ai le plus hésité & \rule{85mm}{0.32pt} \\
Une notion que je veux revoir dès la première semaine & \rule{85mm}{0.32pt} \\
\bottomrule\end{tabularx}
\vfill
\begin{graybox}\small Ce document est un diagnostic de sortie. Il sera analysé compétence par compétence, puis comparé au positionnement passé en début de stage lorsque cette comparaison est légitime.\end{graybox}
\end{document}
